\section{模型假设}
为简化问题、便于建模与求解，本文做出以下合理假设：

\begin{itemize}
  \item 假设 1：……（填写假设，并说明理由）；
  \item 假设 2：……；
  \item 假设 3：……。
\end{itemize}

\section{符号说明}
本文使用的主要符号及其含义如下表所示：

\begin{table}[htbp]
  \centering
  \caption{主要符号说明表}
  \label{tab:symbols}
  \begin{tabularx}{\textwidth}{@{}>{\centering}p{4cm}>{\centering\arraybackslash}X@{}}
    \toprule
    符号 & 含义 \\
    \midrule
    $n$      & 抽样样本量 \\
    $p$      & 次品率 \\
    $X_i$    & 第 $i$ 个样本的观测值 \\
    $\bar{X}$ & 样本均值 \\
    \bottomrule
  \end{tabularx}
\end{table}
